\capitulocapa{images/cap-6.png}{Todo dado pertence ao cliente?}
\label{chap:estado-do-servidor}

\section{A lista que nunca estava realmente pronta}

\begin{historia}
O DevPanel ganhou uma tela de usuários. Na primeira versão, a página fazia uma
requisição ao abrir, guardava o resultado em estado local e exibia uma mensagem
de carregamento.

Pouco depois, a mesma lista apareceu no seletor do modal de convite. A equipe
moveu os usuários para uma store do Zustand e considerou o problema resolvido.
Então surgiram perguntas difíceis: quando buscar novamente? O que acontece se
duas telas solicitarem a lista juntas? Como saber se o conteúdo ainda representa
o servidor? Quem limpa o erro antigo antes de uma nova tentativa?

A lista estava compartilhada, mas continuava sem uma política de sincronização.
\end{historia}

\imagemcapitulo
  {6-1-lista-desatualizada}
  {Cópias consultadas em momentos diferentes deixam a interface inconsistente.}
  {lista-desatualizada}

Compartilhar o resultado de uma requisição não transforma esse resultado em
estado do cliente. Antes de escolher onde guardar um dado, precisamos perguntar
quem é seu verdadeiro proprietário.

\section{Quem pode alterar a fonte de verdade?}

O tema do DevPanel pertence ao cliente. Ele muda quando uma pessoa usa o botão e
a aplicação conhece imediatamente o novo valor. Podemos modelar essa mudança com
uma ação da store.

A lista de usuários pertence ao servidor. Outra pessoa, uma rotina administrativa
ou um serviço externo pode alterá-la sem que a aba aberta seja avisada. A cópia
no navegador é apenas uma fotografia obtida em determinado momento.

Essa diferença produz responsabilidades próprias:

\begin{itemize}
  \item iniciar e repetir requisições;
  \item representar carregamento e falha;
  \item reutilizar respostas já obtidas;
  \item decidir quando uma resposta ficou desatualizada;
  \item sincronizar novamente em momentos apropriados;
  \item evitar requisições idênticas concorrentes.
\end{itemize}

Podemos programar tudo isso em Zustand, mas estaríamos construindo um gerenciador
de consultas dentro de um gerenciador de estado do cliente.

\section{A solução manual e seu custo}

Uma implementação habitual guarda a lista em \codigo{users} e mantém indicadores
separados para carregamento e erro. Uma ação \codigo{fetchUsers} faz a busca; o
componente a chama em um efeito e renderiza conforme essas parcelas.

O código parece pequeno até precisarmos distinguir situações diferentes. A
primeira carga não possui dados para mostrar; uma atualização posterior pode
manter a lista visível enquanto busca uma versão nova. Um erro inicial bloqueia
a tela; um erro durante a atualização pode preservar a resposta anterior.

Também precisamos controlar chamadas duplicadas, descarte de respostas antigas
e tempo de validade. Esses não são detalhes específicos da regra de usuários.
São problemas recorrentes de qualquer consulta remota.

\begin{decisao}
Zustand continua sendo uma boa escolha para estado do cliente compartilhado. Uma
ferramenta de consultas é mais adequada quando o problema central é manter uma
cópia local sincronizada com uma fonte remota.
\end{decisao}

\section{A primeira consulta com TanStack Query}

TanStack Query organiza requisições por chaves e mantém um cache associado a
elas. Primeiro, conectamos um \codigo{QueryClient} à aplicação:

\begin{lstlisting}[caption={Cliente de consultas na raiz da aplicação.}]
import {
  QueryClient,
  QueryClientProvider,
} from '@tanstack/react-query'

const queryClient = new QueryClient()

export function App() {
  return (
    <QueryClientProvider client={queryClient}>
      <AppRoutes />
    </QueryClientProvider>
  )
}
\end{lstlisting}

O Provider não transforma TanStack Query em substituto para Context ou Zustand.
Ele disponibiliza a mesma instância de cache para os componentes consumidores.
O dado continua vindo da API.

A função da requisição fica separada da interface:

\begin{lstlisting}[caption={Função responsável por buscar usuários.}]
export type User = {
  id: string
  name: string
  email: string
}

export async function getUsers(): Promise<User[]> {
  const response = await fetch('/api/users')

  if (!response.ok) {
    throw new Error('Nao foi possivel carregar os usuarios')
  }

  return response.json()
}
\end{lstlisting}

É importante lançar um erro quando a resposta HTTP não for bem-sucedida. O
\codigo{fetch} não rejeita automaticamente a Promise apenas por receber, por
exemplo, um status 500.

\section{Chave, função e estados da consulta}

Na página, \codigo{useQuery} conecta a interface ao cache:

\begin{lstlisting}[caption={Lista de usuários consumindo uma consulta.}]
export function UsersPage() {
  const usersQuery = useQuery({
    queryKey: ['users'],
    queryFn: getUsers,
    staleTime: 60_000,
  })

  if (usersQuery.isPending) return <p>Carregando...</p>
  if (usersQuery.isError) return <p>{usersQuery.error.message}</p>

  return (
    <ul>
      {usersQuery.data.map((user) => (
        <li key={user.id}>{user.name}</li>
      ))}
    </ul>
  )
}
\end{lstlisting}

A chave \codigo{['users']} identifica a consulta. Componentes que usam a mesma
chave observam a mesma entrada do cache. A função explica como obter o dado. Já
\codigo{staleTime} informa por quanto tempo a resposta será considerada fresca;
isso não é o mesmo que apagar o cache depois de um minuto.

\codigo{isPending} representa a ausência de uma primeira resposta. Em uma
atualização posterior, pode haver dados disponíveis enquanto
\codigo{isFetching} indica atividade em segundo plano. Separar essas situações
evita substituir uma lista útil por uma tela vazia a cada sincronização.

\imagemcapitulo
  {6-2-fluxo-query-cache-review-2}
  {Uma entrada do cache atende diferentes consumidores e pode ser revalidada.}
  {fluxo-query-cache}

\section{Cache não significa dado eterno}

Um cache guarda uma resposta para reutilização, mas ela pode ficar obsoleta. O
tempo de frescor expressa quanto tempo aceitamos aquela fotografia sem solicitar
outra automaticamente em determinados eventos.

Os padrões da biblioteca são deliberadamente atentos à sincronização. Consultas
sem um \codigo{staleTime} configurado são consideradas desatualizadas
imediatamente e podem ser buscadas novamente quando um consumidor monta, a janela
recupera o foco ou a rede reconecta. Ajustar esse comportamento exige conhecer a
dinâmica do produto, não escolher um número arbitrário.

Para usuários administrativos, um minuto pode ser razoável. Para permissões
críticas, talvez seja necessário invalidar logo após uma alteração. Para uma
tabela praticamente estática, um período maior reduz trabalho sem comprometer a
experiência.

\section{Zustand e TanStack Query lado a lado}

As duas ferramentas podem coexistir porque resolvem responsabilidades distintas.

\begin{itemize}
  \item Zustand guarda tema, abertura de Drawer e filtros ainda não enviados.
  \item TanStack Query guarda cópias de usuários, relatórios e outros recursos
    obtidos do servidor.
  \item O componente pode combinar as duas fontes sem copiar dados de uma para a
    outra.
\end{itemize}

Evite receber \codigo{data} de uma consulta e copiá-la para uma store por meio de
um efeito. Isso cria duas fontes locais de verdade e exige sincronização manual.
Se a interface precisa transformar a lista, prefira derivar o resultado durante
a leitura ou usar os recursos de seleção da própria consulta.

\imagemcapitulo
  {6-3-cliente-servidor}
  {O DevPanel combina estado do cliente e dados pertencentes ao servidor.}
  {cliente-servidor}

\section{Exercício: atualizar sem duplicar o carregamento}

\begin{tarefa}
Adicione à página de Usuários um botão ``Atualizar lista''. Enquanto a nova
requisição estiver em andamento, mantenha a lista atual visível e altere apenas o
texto do botão para ``Atualizando...''. Não crie outro estado de loading.
\end{tarefa}

Você pode chamar \codigo{refetch} devolvido pela consulta ou invalidar a chave
\codigo{['users']} pelo \codigo{QueryClient}. Registre em seguida:

\begin{enumerate}
  \item quem é proprietário da lista;
  \item qual chave identifica a consulta;
  \item como a primeira carga difere da atualização;
  \item por que a lista não foi copiada para Zustand;
  \item qual evento do produto deveria invalidar esse cache.
\end{enumerate}

\espacoresposta{9}

\begin{resumo}
\begin{itemize}
  \item Estado do servidor é uma cópia local de uma fonte remota.
  \item Uma consulta precisa representar cache, erro, carregamento e
    sincronização.
  \item A chave identifica e permite reutilizar uma entrada do cache.
  \item \codigo{isPending} e \codigo{isFetching} descrevem momentos diferentes.
  \item \codigo{staleTime} define frescor, não a remoção imediata do dado.
  \item Zustand e TanStack Query podem coexistir sem duplicar responsabilidades.
\end{itemize}
\end{resumo}

Escolher uma store começa por classificar a natureza do dado. No próximo
capítulo, voltaremos ao estado do cliente para construir uma store do Zustand a
partir de um contrato pequeno e explícito.
